% ============================================================
% DL01 — From Axioms to Life: The PDL-V Conjecture and the
%         Hierarchical Replication of C1--C4
% Projective Dynamic Logo (PDL) Framework
% Author  : Cédric Laubscher
% Licence : CC BY 4.0
% LaTeX conventions: British English; no spurious mid-sentence
% line breaks; prose as continuous paragraphs;
% \bibliographystyle{unsrt} with natbib [numbers]
% ============================================================

\documentclass[12pt,a4paper]{article}

\usepackage[utf8]{inputenc}
\usepackage[T1]{fontenc}
\usepackage[british]{babel}
\usepackage[numbers]{natbib}
\usepackage{amsmath,amssymb,amsthm}
\usepackage{mathtools}
\usepackage{graphicx}
\usepackage{float}
\usepackage{booktabs}
\usepackage{longtable}

\usepackage{array}
\usepackage{xcolor}
\usepackage{hyperref}
\usepackage{geometry}
\usepackage{enumitem}
\usepackage[protrusion=true,expansion=true,tracking=true,kerning=true,spacing=true,activate={true,nocompatibility},final,stretch=10,shrink=10,step=2]{microtype}
\microtypecontext{spacing=nonfrench}

\usepackage{setspace}
\onehalfspacing
\usepackage{tcolorbox}
\tcbuselibrary{theorems,skins,breakable}

\geometry{top=2.5cm,bottom=2.5cm,left=2.8cm,right=2.8cm}

\hypersetup{
  colorlinks=true,
  linkcolor=blue!70!black,
  citecolor=blue!60!black,
  urlcolor=blue!70!black,
  pdftitle={From Axioms to Life: The PDL-V Conjecture and the
            Hierarchical Replication of C1--C4},
  pdfauthor={Cédric Laubscher}}

% ---- Colour palette ----------------------------------------
\definecolor{proofblue}{RGB}{0,70,140}
\definecolor{conjecturegreen}{RGB}{0,120,60}
\definecolor{openproblemred}{RGB}{160,20,20}
\definecolor{resolvedgreen}{RGB}{0,130,70}
\definecolor{chaincolor}{RGB}{30,100,180}
\definecolor{epscolor}{RGB}{160,60,0}
\definecolor{darkgray}{RGB}{80,80,80}
\definecolor{lifeblue}{RGB}{0,90,160}
\definecolor{consciousviolet}{RGB}{100,0,140}

% ---- Theorem environments ----------------------------------
\theoremstyle{plain}
\newtheorem{theorem}{Theorem}[section]
\newtheorem{lemma}[theorem]{Lemma}
\newtheorem{proposition}[theorem]{Proposition}
\newtheorem{corollary}[theorem]{Corollary}

\theoremstyle{definition}
\newtheorem{definition}[theorem]{Definition}
\newtheorem{axiom}{Axiom}

\theoremstyle{remark}
\newtheorem{remark}[theorem]{Remark}

% ---- Custom tcolorbox environments -------------------------
\tcbset{
  theoremstyle/.style={colframe=proofblue,colback=blue!3,
    fonttitle=\bfseries\color{proofblue},breakable,enhanced},
  conjecturestyle/.style={colframe=conjecturegreen,colback=green!3,
    fonttitle=\bfseries\color{conjecturegreen},breakable,enhanced},
  openproblemstyle/.style={colframe=openproblemred,colback=red!3,
    fonttitle=\bfseries\color{openproblemred},breakable,enhanced},
  resolvedstyle/.style={colframe=resolvedgreen,colback=green!4,
    fonttitle=\bfseries\color{resolvedgreen},breakable,enhanced},
  computedstyle/.style={colframe=chaincolor,colback=chaincolor!4,
    fonttitle=\bfseries\color{chaincolor},breakable,enhanced},
  lifestyle/.style={colframe=lifeblue,colback=lifeblue!4,
    fonttitle=\bfseries\color{lifeblue},breakable,enhanced}}

\newenvironment{theorembox}[1][]{%
  \begin{tcolorbox}[theoremstyle,title={Theorem\ifx\relax#1\relax\else: #1\fi}]}{\end{tcolorbox}}
\newenvironment{conjecturebox}[1][]{%
  \begin{tcolorbox}[conjecturestyle,title={Conjecture\ifx\relax#1\relax\else: #1\fi}]}{\end{tcolorbox}}
\newenvironment{openproblem}[1][]{%
  \begin{tcolorbox}[openproblemstyle,title={Open Problem\ifx\relax#1\relax\else: #1\fi}]}{\end{tcolorbox}}
\newenvironment{computedresult}[1][]{%
  \begin{tcolorbox}[computedstyle,title={Computational Result\ifx\relax#1\relax\else: #1\fi}]}{\end{tcolorbox}}
\newenvironment{lifebox}[1][]{%
  \begin{tcolorbox}[lifestyle,title={#1}]}{\end{tcolorbox}}

% ---- Macros ------------------------------------------------
\newcommand{\PDL}{\textsc{PDL}}
\newcommand{\Kfour}{K_4}
\newcommand{\Calg}{\mathcal{C}}
\newcommand{\Sop}{\mathcal{S}}
\newcommand{\Rop}{\mathcal{R}}
\newcommand{\Lop}{\mathcal{L}}
\newcommand{\eps}{\varepsilon}
\newcommand{\deltastar}{\delta^{*}}

% ============================================================
\begin{document}
% ============================================================

\title{\textbf{From Axioms to Life:\\
The \PDL-V Conjecture and the\\
Hierarchical Replication of C1--C4}\\[0.5em]
{\large \PDL\ Framework --- Document DL01}}

\author{Cédric Laubscher\\
\small Independent Researcher, Switzerland\\
\small \href{https://cedriclaubscher.ch}{cedriclaubscher.ch}}

\date{May 2026}
\maketitle

% ============================================================
\begin{abstract}
The Projective Dynamic Logo (\PDL) framework derives the fundamental constants of
physics --- the proton mass, Newton's constant, the cosmological constant --- from
four combinatorial axioms C1--C4 on finite signed graphs, without presupposing
spacetime, particles, or fields. Documents D01--D53 establish this derivation as a
closed causal chain with no free parameter beyond the quark mass difference
$\Delta m_{\mathrm{iso}}$. The present document opens the DL series, a new programme
that asks whether the same four axioms, applied hierarchically to composite closures,
force the emergence of life and consciousness as theorems rather than as mysteries.
We introduce the lifting operator $\Lop$, the constraint-structure operator $\Sop(\Gamma)$,
and the self-representation relation $\Rop^k(\Gamma)$. We state the \PDL-V conjecture
in four parts: the existence of a life threshold $n^{*}_{\mathrm{life}}$, its selection
by C4, the necessity of the birth--death cycle, and the existence of a thermodynamic
ceiling $n^{*}_{\mathrm{max}}$ forced by combinatorial entropy. Five computational
scripts verify the foundational claims on small cases: $\Sop(K_4) \cong K_4$
unconditionally (8/8 coherent configurations), weak self-representation
$\Rop^{1}_{\mathrm{weak}}$ is forced at 100\% for all $n \geq 4$, and the rarity law
$f(n) = 2^{n+1-\binom{n}{2}}$ establishes $n^{*}_{\mathrm{max}}$ as a finite combinatorial
necessity. The synaptic mechanism is presented as the first mechanistic illustration
of the lifting operator at the neural level. All claims are stratified by epistemic
status; the \PDL-V conjecture is not yet a theorem.
\end{abstract}

\tableofcontents
\newpage

% ============================================================
\section{Why This Document Exists}
\label{sec:motivation}
% ============================================================

The \PDL\ framework was constructed as a foundational programme in physics. Its axioms
C1--C4 are stated at the level of finite signed graphs, prior to any presupposition of
space, time, matter, or energy. From these four axioms alone, documents D01--D53
derive the electron Compton oscillation (D01--D02), the proton quintuplet
$(24,28,930,10087,11017)$ (D05--D10), Newton's constant $G$ to 27\,ppm (D36--D44),
the Dirac and Schrödinger equations (D32--D33), Born's rule (D34, D46), the periodic
table of elements (D47), the London equation of superconductivity (D49), the
Bekenstein--Hawking coefficient $1/4$ (D50), and the cosmological constant $\Lambda$
(D51--D52). The causal chain C1--C4 $\to$ $G$ is closed, with no free parameter
beyond $\Delta m_{\mathrm{iso}} = m_d - m_u$, the quark mass difference at the
\PDL--QCD interface \cite{DM_v18}.

A natural question then presents itself with a force that is difficult to ignore. If
C1--C4 are prior to physics --- if physics is their first theorem --- are they also
prior to life and consciousness? If the same four axioms, applied iteratively to
composite closures, force the emergence of self-replicating structures and eventually
of self-representing ones, then life and consciousness are not mysteries layered on
top of physics. They are theorems in the same causal chain, separated from the proton
only by the number of hierarchical levels traversed.

A note on context is essential before proceeding. The reader unfamiliar with the
\PDL\ foundation is strongly encouraged to consult the Global Mapping document
DM~v18 \cite{DM_v18} before reading this document. The DL programme is an extension
of an established foundational programme that has derived numerical predictions
verified to parts per million --- it is not a standalone speculative proposal.
The present document should be read in that light: its conjectures are grounded in
a formally closed physical foundation, not in intuition alone.

This document does not claim to have proved that life or consciousness are theorems
of C1--C4. What it does is the following. It introduces the formal objects necessary
to make the question precise, presents the \PDL-V conjecture as a structured research
programme, reports the results of five computational verification scripts on small
cases, and situates the programme with respect to existing literature. Every claim is
labelled with its epistemic status. The reader who disagrees is invited to identify
which axiom, which definition, or which computational result is incorrect. That is the
only form of engagement the programme solicits.

\subsection{Epistemic stratification}

Throughout this document and the DL series, the following classification is used
without exception.

\begin{center}
\begin{tabular}{p{4.1cm}p{10.4cm}}
\toprule
\textbf{Status} & \textbf{Meaning} \\
\midrule
Unconditional theorem & Derived from C1--C4 alone, verified computationally \\
Conditional theorem & Derived under named hypotheses \\
Computational result & Verified exhaustively by script on finite cases \\
Structured conjecture & Formally stated, not yet proved, with falsification criteria \\
Interpretive remark & Coherent with \PDL, not derived \\
Open problem & Identified gap requiring future work \\
\bottomrule
\end{tabular}
\end{center}

% ============================================================
\section{The PDL Foundation: A Self-Contained Summary}
\label{sec:foundation}
% ============================================================

This section is self-contained. A reader with access to D01--D10 and D29 may read it
as a condensed reference; a reader new to \PDL\ will find here everything needed to
follow the DL programme.

\subsection{The four axioms}

The \PDL\ framework is built on four axioms stated on finite signed graphs. A
\emph{signed graph} is a pair $(G, s)$ where $G = (V, E)$ is a finite graph and
$s: E \to \{+1, -1\}$ assigns a sign to each edge.

\begin{axiom}[C1 --- Binary pulsation]
The configuration admits a logical 2-cycle: a sign assignment that returns to itself
after exactly two update steps and is not a fixed point. Existence is the minimal
form of repetition.
\end{axiom}

\begin{axiom}[C2 --- Triangular coherence]
Every triangle (closed triple of edges) satisfies $s_{ij} \cdot s_{jk} \cdot s_{ik} = +1$.
A triangle with product $+1$ is \emph{coherent}; one with product $-1$ is
\emph{frustrated}. The coherence cost $\eps(\Gamma)$ is the fraction of frustrated
triangles.
\end{axiom}

\begin{axiom}[C3 --- Minimal completeness]
The configuration is irreducible: it admits no non-trivial decomposition into
independent sub-configurations each satisfying C1 and C2 separately. The underlying
graph is connected and its coherence structure is non-reducible.
\end{axiom}

\begin{axiom}[C4 --- Logical optimisation]
Among all configurations satisfying C1--C3, those that are physically instantiated
are those that minimise the coherence cost $\eps(\Gamma)$ for a given number of
vertices and edges.
\end{axiom}

\subsection{The minimal admissible closure \texorpdfstring{$K_4$}{K4}}

\begin{theorem}[D01--D02, \cite{D01,D02}]
No finite signed graph on fewer than four vertices satisfies all four axioms
simultaneously. The complete signed graph on four vertices $V = \{0,1,2,3\}$ and six
edges, denoted $\Kfour$, is the first admissible elementary stationary closure. It is
minimal (no proper sub-closure satisfies C1--C4) and quasi-unique (the sign pattern
is unique up to the global symmetry group $S_4$).
\end{theorem}

$\Kfour$ undergoes a binary pulsation alternating between a coherent configuration
$s^{(1)}$ and its global flip $s^{(2)} = -s^{(1)}$. There are exactly eight coherent
sign configurations on $\Kfour$, forming four pulsation pairs $\{s, -s\}$. The
coherence cost of any coherent configuration is $\eps = 0$ by definition. This
pulsation is identified with the Compton oscillation of the electron; its minimal
quantum of action is $\hbar$ \cite{D01}.

\subsection{The \texorpdfstring{$(A)\wedge(B)$}{(A)∧(B)} coupling criterion}

The coupling of two closures is governed by a stability criterion derived in D29
\cite{D29}. A \emph{mixed triangle} couples two $\Kfour$ vertices to one proton
vertex. It is $(A)\wedge(B)$-stable if and only if its cross-product satisfies both
the alignment condition (A) and the exact sign inversion between the two half-cycles
(B). Exhaustive verification over 768 configurations yields 192 stable ones. The
ratio $192/768 = 1/4$ is the fundamental coupling probability $\kappa$, which is not
a free parameter but a combinatorial consequence of C1--C4.

\subsection{What the corpus has established}

The following results are unconditional theorems of C1--C4, verified computationally
and published on Zenodo \cite{DM_v18}. The proton mass $m_p$ and proton--electron
mass ratio $\mu^* = 1836.152670$ are derived at 27\,ppm (D05--D10, D30). Newton's
constant $G$ is derived with no free parameter (D36--D44). The Dirac equation, the
Schrödinger equation, and Born's rule emerge from the $(A)\wedge(B)$ coupling
criterion applied to the $\Kfour$ pulsation (D32--D34, D46). The periodic table, the
magic numbers of nuclear physics, and the shell-filling rates $r_{\mathrm{exc}}(Z)$
for $Z = 1$--$82$ are unconditional theorems (D47). The London equation of
superconductivity and the Bekenstein--Hawking coefficient $1/4$ are theorems (D49,
D50). The cosmological constant $\Lambda$ is derived without free parameter
(D51--D52).

This enumeration is not rhetorical. It establishes the epistemic authority of the
foundation from which the DL programme departs.

% ============================================================
\section{The Scale-Invariance of C1--C4}
\label{sec:scaleinvariance}
% ============================================================

We have recalled the foundation. We now take the first step beyond it. The key
observation is simple to state and far-reaching in its consequences: the four axioms
C1--C4 do not belong to the quantum level. They belong to existence at any scale.
This section makes that claim precise, and states it as a formal conjecture that the
rest of the DL programme will work to prove. The axioms C1--C4
are not axioms \emph{of the quantum level}. They are axioms \emph{of existence at
any scale}. This section makes that claim precise.

\subsection{C1 at every scale}

Axiom C1 does not say that $\Kfour$ pulsates. It says that any configuration that
exists must admit a 2-cycle. Existence and cyclic repetition are identified. The
consequences of this identification propagate upward through every level of the
hierarchy without modification.

Consider the following sequence of 2-cycles observed at increasing scales. At the
quantum level, $\Kfour$ alternates between $s^{(1)}$ and $s^{(2)} = -s^{(1)}$ at
the Compton frequency. At the atomic level, electrons transition between energy
levels in discrete cycles. At the molecular level, chemical bonds vibrate and rotate
in periodic modes. At the cellular level, the cell cycle --- growth, replication,
division --- is a 2-cycle in the logical sense: the daughter cell is the configuration
returned after one complete update of the parent. At the organismal level, the
heartbeat, the respiratory cycle, and the sleep--wake cycle are all expressions of
C1 at their respective timescales. At the lineage level, birth and death constitute
the 2-cycle of the generation: the offspring is the configuration returned after one
complete generational update of the parent.

None of these cycles are imposed from outside the framework. They are the expression
at each scale of the single condition that makes existence possible: a configuration
must return to a form of itself, or it dissolves.

\begin{remark}
The brick broken by a hammer illustrates C1 at the geological scale. The fracture
does not create new matter; it reveals sub-closures that were already coherent
internally (C3) and that become autonomous when the constraints holding them within
the larger closure are released. The fragments are not copies of the brick; they are
configurations that C4 had already selected as locally optimal within the brick's
structure.
\end{remark}

\subsection{The conjecture of scale-invariance}

The observation above can be stated as a formal conjecture.

\begin{conjecturebox}[Scale-invariance of C1--C4]
Let $\Gamma \in \Calg_n$ be an admissible closure of level $n$, and let
$\Gamma' = \Lop(\Gamma_1, \ldots, \Gamma_k)$ be a composite closure of level $n+1$
constructed by the lifting operator $\Lop$ (Definition~\ref{def:lifting}). Then
$\Gamma'$ satisfies C1--C4 with renormalised parameters. In particular, the
coherence cost $\eps(\Gamma')$, the coupling probability $\kappa'$, and the pulsation
period $T'$ are determined by the parameters of the level-$n$ closures via
$\Lop$, without invoking any new axiom.
\end{conjecturebox}

\emph{Status: structured conjecture. Partially verified computationally for $n = 0
\to 1$ (Scripts 1--5). Not yet proved analytically.}

% ============================================================
\section{The Lifting Operator \texorpdfstring{$\Lop$}{L}}
\label{sec:lifting}
% ============================================================

\subsection{Intuition before the definition}

Before stating the formal definition, consider what happens when four protons come
close enough to interact. Each proton is itself a composite closure built from $\Kfour$
instances via the $(A)\wedge(B)$ coupling criterion. At the nuclear level, each proton
becomes a \emph{point} --- an indivisible object at that scale, exactly as $\Kfour$
is a point at the quantum level. The four protons then form a new signed graph: they
are the vertices, the nuclear forces between them are the edges, and the signs of
those edges are determined by whether the pulsations of the two protons are in phase
($+1$) or in opposition ($-1$). C1--C4 are then applied to this new graph, selecting
the configurations of minimal coherence cost. The result is the helium-4 nucleus.

This is precisely what the lifting operator does, at every level of the hierarchy.

\subsection{Definition}

\begin{definition}[Lifting operator $\Lop$]
\label{def:lifting}
Let $\Calg_n$ denote the set of admissible closures of level $n$. The lifting
operator is a map
\begin{equation}
\Lop : \mathcal{P}(\Calg_n) \to \Calg_{n+1}
\end{equation}
defined as follows. Given a finite set $\{\Gamma_1, \ldots, \Gamma_k\} \subset
\Calg_n$ with $k \geq 4$, the composite closure $G_{n+1} = \Lop(\Gamma_1, \ldots,
\Gamma_k)$ is the signed graph whose vertex set is $\{\Gamma_1, \ldots, \Gamma_k\}$,
whose edges are determined by the $(A)\wedge(B)$ coupling criterion between pairs of
level-$n$ closures, and whose edge signs are determined by the phase parity of the
coupling: $s_{ij} = +1$ if the pulsations of $\Gamma_i$ and $\Gamma_j$ are in phase
at the first half-cycle, and $s_{ij} = -1$ if they are in opposition. The admissible
configuration of $G_{n+1}$ is selected by C4 among all configurations compatible
with C2.
\end{definition}

The minimum $k \geq 4$ is not arbitrary: it is forced by C1 applied to $G_{n+1}$.
For $G_{n+1}$ to admit a 2-cycle, the same combinatorial argument that forces
$\Kfour$ at level~0 forces at least four vertices at every subsequent level. A
composite closure built from fewer than four sub-closures cannot satisfy C1 and is
therefore not admissible.

\subsection{The proton as the first application of \texorpdfstring{$\Lop$}{L}}

The construction of the proton from $\Kfour$ in documents D05--D10 is precisely the
first application of $\Lop$. The proton is $\Lop(\Kfour^{\times k})$ for an
appropriate $k$, where the coupling is governed by the $(A)\wedge(B)$ criterion. The
proton quintuplet $(24, 28, 930, 10087, 11017)$ is the output of C4 applied to this
construction. The lifting operator is therefore not a new object introduced for the
DL programme; it is the mechanism already operative in the foundational corpus,
named and formalised here for application at higher levels.

\subsection{Self-similarity of \texorpdfstring{$\Lop$}{L}}

A key property observed computationally (Script 1) is that the level-1 closure
produced by $\Lop$ from four $\Kfour$ instances has the same structural form as
$\Kfour$ itself: four vertices, six edges, four triangles. This self-similarity is
the formal counterpart of the scale-invariance conjecture. It says that the lifting
operator does not produce qualitatively new structures at each level --- it produces
the same structure, populated by objects of increasing complexity.

% ============================================================
\section{The Constraint-Structure Operator
         \texorpdfstring{$\Sop(\Gamma)$}{S(Γ)}}
\label{sec:S}
% ============================================================

The lifting operator $\Lop$ tells us how closures combine into composites. We now
need a second operator that asks a different question: does a closure \emph{know}
its own structure? Can it encode, internally, the rules by which it holds together?
The constraint-structure operator $\Sop$ formalises this question precisely, and its
answer for $\Kfour$ will turn out to be the most important single result of this
document.

\subsection{Definition}

\begin{definition}[Constraint-structure operator $\Sop$]
\label{def:S}
Let $\Gamma \in \Calg_n$ be an admissible closure with vertex set $V$, edge set $E$,
and sign assignment $s: E \to \{+1,-1\}$. The \emph{constraint-structure graph}
$\Sop(\Gamma)$ is the signed graph defined as follows. Its vertex set is the set of
triangles of $\Gamma$: each triangle $\tau_{ijk}$ becomes a vertex of $\Sop(\Gamma)$.
Two vertices $\tau_{ijk}$ and $\tau_{ijl}$ of $\Sop(\Gamma)$ are connected by an
edge if and only if the corresponding triangles share an edge in $\Gamma$. The sign
of that edge in $\Sop(\Gamma)$ is the sign of the shared edge in $\Gamma$.
\end{definition}

The operator $\Sop$ maps a signed graph to another signed graph of the same formal
type. It encodes the dependency structure of the coherence constraints: two
constraints are adjacent in $\Sop(\Gamma)$ if they share a common edge, and the sign
of their adjacency is the sign of the edge they share.

\subsection{Worked example: \texorpdfstring{$\Sop(\Kfour)$}{S(K4)} step by step}

To make Definition~\ref{def:S} concrete, we compute $\Sop(\Kfour)$ explicitly for
the coherent configuration $s = (-1,-1,+1,+1,-1,-1)$ on the six edges
$(e_{01}, e_{02}, e_{03}, e_{12}, e_{13}, e_{23})$.

\emph{Step 1 --- Identify the triangles of $\Kfour$.} There are $\binom{4}{3} = 4$
triangles, which become the four vertices of $\Sop(\Kfour)$:
\begin{align*}
v_0 &= \tau_{012} = \text{triangle on vertices } \{0,1,2\}, \\
v_1 &= \tau_{013} = \text{triangle on vertices } \{0,1,3\}, \\
v_2 &= \tau_{023} = \text{triangle on vertices } \{0,2,3\}, \\
v_3 &= \tau_{123} = \text{triangle on vertices } \{1,2,3\}.
\end{align*}

\emph{Step 2 --- Determine the edges of $\Sop(\Kfour)$.} Two triangles are adjacent
in $\Sop(\Kfour)$ if they share an edge in $\Kfour$. Every pair of triangles in
$\Kfour$ shares exactly one edge (there are $\binom{4}{2} = 6$ pairs, each sharing
one of the six edges of $\Kfour$). The six edges of $\Sop(\Kfour)$ and their
inherited signs are:
\begin{center}
\begin{tabular}{llc}
\toprule
Edge in $\Sop(\Kfour)$ & Shared edge in $\Kfour$ & Sign \\
\midrule
$(v_0, v_1)$ & $e_{01}$ & $-1$ \\
$(v_0, v_2)$ & $e_{02}$ & $-1$ \\
$(v_0, v_3)$ & $e_{12}$ & $+1$ \\
$(v_1, v_2)$ & $e_{03}$ & $+1$ \\
$(v_1, v_3)$ & $e_{13}$ & $-1$ \\
$(v_2, v_3)$ & $e_{23}$ & $-1$ \\
\bottomrule
\end{tabular}
\end{center}

\emph{Step 3 --- Conclusion.} $\Sop(\Kfour)$ has 4 vertices and 6 edges --- the same
counts as $\Kfour$ itself. It is structurally isomorphic to $\Kfour$. This holds for
all eight coherent configurations of $\Kfour$, as verified exhaustively by Script~1.

\subsection{The fixed-point theorem for \texorpdfstring{$K_4$}{K4}}

\begin{computedresult}[$\Sop(\Kfour) \cong \Kfour$ --- Script~1]
\label{thm:fixedpoint}
For every one of the eight coherent sign configurations of $\Kfour$, the
constraint-structure graph $\Sop(\Kfour)$ has exactly four vertices and six edges,
and is structurally isomorphic to $\Kfour$. This result holds unconditionally: it
is independent of the choice of coherent configuration.
\end{computedresult}

\begin{proof}[Verification]
$\Kfour$ has four triangles: $(0,1,2)$, $(0,1,3)$, $(0,2,3)$, $(1,2,3)$. These
become the four vertices of $\Sop(\Kfour)$. Every pair of triangles in $\Kfour$
shares exactly one edge (since $\Kfour$ is the complete graph on four vertices, every
pair of its four triangular faces has a common edge). Therefore $\Sop(\Kfour)$ has
$\binom{4}{2} = 6$ edges. A signed graph on four vertices with six edges is $\Kfour$.
Exhaustive verification over all eight coherent configurations confirms the result
(Script~1, 8/8).
\end{proof}

This theorem says that $\Kfour$ is a fixed point of $\Sop$: the graph of its
coherence constraints is isomorphic to itself. This is a structural property of
$\Kfour$ alone: for $n \geq 5$, the number of triangles in $K_n$ is $\binom{n}{3}
> n$, so $\Sop(K_n)$ cannot be isomorphic to $K_n$ for $n \geq 5$ (Script~1,
confirmed 0/16 for $n=5$ and 0/32 for $n=6$).

\begin{remark}
The fixed-point property of $\Kfour$ under $\Sop$ is not a coincidence. It reflects
the fact that $\Kfour$ is the unique finite graph whose triangles are in bijection
with its vertices and whose inter-triangle adjacencies reproduce its edge structure.
This is why $\Kfour$ is the seed of the entire hierarchy: it is the only level at
which the self-referential structure of $\Sop$ closes exactly on itself.
\end{remark}

% ============================================================
\section{The Self-Representation Relation
         \texorpdfstring{$\Rop^k(\Gamma)$}{Rk(Γ)}}
\label{sec:R}
% ============================================================

We now have the tools to formalise the central concept of the DL programme. A closure
can carry within itself an encoding of its own constraints --- a description of the
rules by which it coheres. This is self-representation. We define it in two versions
--- strong and weak --- because the distinction turns out to matter enormously: it
is the difference between a crystal and a living cell.

\subsection{Two definitions of self-representation}

The concept of self-representation has two natural formalisations in this framework,
of different strength.

\begin{definition}[Strong self-representation $\Rop^1_{\mathrm{strong}}$]
A closure $\Gamma \in \Calg_n$ is \emph{strongly self-representing} if it contains a
sub-graph $\Gamma' \subseteq \Gamma$ such that $\Gamma' \cong \Sop(\Gamma)$ as signed
graphs, and $\Sop(\Gamma)$ is itself coherent (satisfies C2).
\end{definition}

\begin{definition}[Weak self-representation $\Rop^1_{\mathrm{weak}}$]
A closure $\Gamma \in \Calg_n$ is \emph{weakly self-representing} if there exists a
sub-graph $\Gamma' \subseteq \Gamma$ and a sign-preserving graph homomorphism
$\varphi: \Gamma' \to \Sop(\Gamma)$. That is, $\Gamma'$ maps into $\Sop(\Gamma)$
with preservation of edge signs, without requiring global isomorphism.
\end{definition}

The strong definition requires an exact internal copy of the constraint structure.
The weak definition requires only that the closure contains a faithful encoding of
its own constraints --- a compressed but accurate representation, not a perfect
replica. The distinction matters: DNA does not contain an isomorphic copy of the
organism; it contains a compressed encoding of the rules by which the organism
constructs itself.

\subsection{Computational results on self-representation}

\begin{computedresult}[Strong self-representation --- Script~1 and 4]
Strong self-representation $\Rop^1_{\mathrm{strong}}$ is present in exactly 4 of
the 8 coherent configurations of $\Kfour$ (those for which $\Sop(\Kfour)$ is
itself coherent). For $n \geq 5$, strong self-representation is absent from all
coherent configurations of $K_n$ (0/16 for $n=5$, 0/32 for $n=6$, 0/64 for $n=7$,
0/128 for $n=8$).
\end{computedresult}

\begin{computedresult}[Weak self-representation --- Script~5]
\label{thm:R1weak}
Weak self-representation $\Rop^1_{\mathrm{weak}}$ is forced at 100\% for every
coherent configuration at $n \geq 4$, and is absent at $n = 3$. Specifically:
$n=3$: 0/4; $n=4$: 8/8; $n=5$: 16/16; $n=6$: 32/32; $n=7$: 64/64.
\end{computedresult}

This result has a precise interpretation. Every closure that satisfies C1--C4 with
four or more vertices automatically contains a faithful encoding of its own coherence
constraints, accessible via a sign-preserving homomorphism. Self-representation in
the weak sense is not an additional property acquired by some closures: it is a
structural consequence of coherence for $n \geq 4$.

\subsection{The hierarchy of self-representation}

The hierarchy of self-representation levels is defined by iterating $\Sop$.

\begin{definition}[Hierarchy $\Rop^k$]
For $k \geq 1$, a closure $\Gamma$ satisfies $\Rop^k_{\mathrm{weak}}$ if there
exists a sub-graph $\Gamma^{(k)} \subseteq \Gamma$ and a sign-preserving homomorphism
$\varphi: \Gamma^{(k)} \to \Sop^k(\Gamma)$, where $\Sop^k$ denotes the $k$-fold
application of $\Sop$.
\end{definition}

The sequence $\Rop^1 \subset \Rop^2 \subset \cdots$ defines a tower of increasingly
reflexive self-represen-tation levels. $\Rop^1$ corresponds to a closure that encodes
its own constraints; $\Rop^2$ corresponds to a closure that encodes the encoding of
its own constraints; and so on. The DL programme conjectures that life corresponds
to $\Rop^1_{\mathrm{active}}$ and consciousness to $\Rop^2_{\mathrm{active}}$, where
\emph{active} denotes the additional condition that the encoding sub-graph
$\Gamma^{(k)}$ is coupled via $(A)\wedge(B)$ to the production mechanism of the
closure.

% ============================================================
\section{The PDL-V Conjecture}
\label{sec:PDLV}
% ============================================================

The \PDL-V conjecture is stated in four parts. Each part is labelled with its
epistemic status.

\begin{conjecturebox}[PDL-V, Part 1 --- Existence of the life threshold]
There exists a minimal integer $n^{*}_{\mathrm{life}} \geq 1$ such that for all
$n \geq n^{*}_{\mathrm{life}}$, every admissible closure $\Gamma \in \Calg_n$
selected by C4 satisfies $\Rop^1_{\mathrm{active}}$: it contains a sub-graph
$\Gamma'$ that encodes $\Sop(\Gamma)$ via a sign-preserving homomorphism, and
$\Gamma'$ is coupled via $(A)\wedge(B)$ to a production mechanism that generates
new closures of the same type with controlled variation.
\end{conjecturebox}

\emph{Status: structured conjecture. Computational results establish that
$\Rop^1_{\mathrm{weak}}$ (the structural precondition) is universally present for
$n \geq 4$. The active coupling component remains to be verified.}

\begin{conjecturebox}[PDL-V, Part 2 --- Selection by C4]
Closures satisfying $\Rop^1_{\mathrm{active}}$ have a strictly lower lineage
coherence cost $\bar{\eps}(\mathcal{L}) = \lim_{N \to \infty} \frac{1}{N}
\sum_{i=0}^{N} \eps(\Gamma_i)$ than closures that persist without active
self-representation. C4 applied at the lineage level therefore selects
$\Rop^1_{\mathrm{active}}$ closures systematically over persistent non-representing
ones.
\end{conjecturebox}

\emph{Status: structured conjecture. The intuitive argument is that a closure which
encodes its own constraints can anticipate coherence violations and correct them
before they propagate, reducing the average coherence cost over the lineage. This
advantage is not yet proved formally.}

\begin{conjecturebox}[PDL-V, Part 3 --- Necessity of the birth--death cycle]
Lineages of finite-duration closures with controlled inter-generational variation
$\delta = \delta^{*}$ have a strictly lower lineage coherence cost than immortal
persistent closures at the same level. Death is a logical necessity imposed by C4
at the lineage scale: an immortal closure accumulates coherence violations without
the possibility of reset, whereas a lineage corrects its coherence structure at each
generational transition. The optimal variation $\delta^{*} = 1/2$ is forced by C4
and verified computationally.
\end{conjecturebox}

\begin{computedresult}[Optimal variation $\delta^{*}$ --- Script~3]
The optimal Hamming distance between successive coherent configurations of a
level-1 closure is $\delta^{*} = 3/6 = 1/2$: exactly half of the six edges change
sign between optimal successive configurations. This value coincides with the
normalised distance between the two half-cycles of the $\Kfour$ pulsation (full
sign inversion = distance 1; half-cycle transition = distance 1/2).
\end{computedresult}

\begin{conjecturebox}[PDL-V, Part 4 --- Thermodynamic ceiling $n^{*}_{\mathrm{max}}$]
There exists a maximal integer $n^{*}_{\mathrm{max}} < \infty$, determined by the
balance between the coherence cost reduction gained by each additional level of
self-representation and the combinatorial entropy cost of maintaining coherence at
that level, such that closures at level $n > n^{*}_{\mathrm{max}}$ are
thermodynamically unstable and fragment towards lower levels. The sequence
$n^{*}_{\mathrm{life}} < n^{*}_{\mathrm{consciousness}} < \cdots <
n^{*}_{\mathrm{max}}$ is finite and forced by C1--C4.
\end{conjecturebox}

\begin{computedresult}[Rarity law and $n^{*}_{\mathrm{max}}$ --- Script~3]
\label{thm:rarity}
The fraction of coherent configurations among all signed configurations on $K_n$
satisfies the exact law
\begin{equation}
f(n) = 2^{n + 1 - \binom{n}{2}}.
\end{equation}
The sequence of exponents $n + 1 - \binom{n}{2}$ for $n = 3, 4, 5, 6$ is
$-1, -3, -6, -10$: the negatives of the triangular numbers $T_{n-2}$. The fraction
$f(n)$ decreases doubly exponentially with $n$, establishing that there exists a
finite $n^{*}_{\mathrm{max}}$ beyond which no coherent closure can be
thermodynamically maintained against fluctuations.
\end{computedresult}

\subsection{Falsification criteria}

The \PDL-V conjecture is falsified by any one of the following.

\begin{enumerate}[leftmargin=2em]
\item Construction of an admissible closure at level $n \geq n^{*}_{\mathrm{life}}$
      that satisfies C1--C4 but does not satisfy $\Rop^1_{\mathrm{active}}$, despite
      C4 selecting it as a coherence minimum.
\item A proof that the lineage coherence cost $\bar{\eps}(\mathcal{L})$ is not
      minimised by $\Rop^1_{\mathrm{active}}$ closures, i.e., that immortal closures
      have equal or lower lineage cost.
\item A proof that $\delta^{*} \neq 1/2$ at any level above the first, i.e., that
      the optimal inter-generational variation is not the half-pulsation of $\Kfour$.
\item A proof that $f(n)$ does not decrease doubly exponentially, i.e., that the
      rarity law is not exact.
\end{enumerate}

% ============================================================
\section{A Concrete Bridge: The Synaptic Mechanism}
\label{sec:synapse}
% ============================================================

We have now introduced the formal objects --- $\Lop$, $\Sop$, $\Rop^k$ --- and stated
the \PDL-V conjecture. Before proceeding to the thermodynamic analysis and the
literature review, it is worth pausing to show that the formalism is not ad hoc. The
synaptic mechanism of the neural network is a concrete, mechanistically well-understood
biological system that illustrates every element of the framework developed above.
This section is an illustration, not a proof. Every claim is labelled accordingly.
We emphasise explicitly that this section draws structural parallels, not formal
derivations. The derivation of Hebbian plasticity from C4, and of the STDP rule from
the half-cycle structure of C1, are the programme of DL04 --- they are not claims
of the present document. What is claimed here is only that the \PDL\ formalism and
the neuroscientific mechanisms are structurally coherent, which motivates the formal
programme without prejudging its outcome.

\subsection{The minimal neural circuit as a signed triangle}

A chemical synapse connects two neurons across a junction of approximately 20\,nm.
Neurotransmitters released from the presynaptic terminal bind to receptors on the
postsynaptic membrane and produce one of two effects: excitation ($s = +1$ on the
edge between the two neurons in the neural graph) or inhibition ($s = -1$). The
action potential itself --- the all-or-nothing electrical pulse --- is a binary
pulsation in the sense of C1: depolarisation followed by repolarisation is a 2-cycle
at the cellular timescale.

Consider the simplest non-trivial neural circuit: three neurons $A$, $B$, $C$
forming a triangle, where $A$ excites $B$ ($s_{AB} = +1$), $B$ inhibits $C$
($s_{BC} = -1$), and $C$ excites $A$ ($s_{CA} = +1$). The coherence of this triangle
in the sense of C2 is:
\begin{equation}
s_{AB} \cdot s_{BC} \cdot s_{CA} = (+1)(-1)(+1) = -1.
\end{equation}
This is a \emph{frustrated} triangle --- a violation of C2. In the \PDL\ framework,
a frustrated triangle under C2 is not a stable configuration: C4 drives the system
towards configurations that reduce the fraction of frustrated triangles. The result
in this circuit is oscillation: the three neurons fire in sequence, never reaching a
static equilibrium, because no static assignment of firing states can satisfy the
coherence condition. This is precisely the behaviour observed in biological central
pattern generators --- the neural circuits that produce rhythmic motor patterns such
as breathing, walking, and heartbeat.

\begin{remark}
The three-neuron frustrated circuit is the neural-level counterpart of the frustrated
triangle at the $\Kfour$ level. At both scales, C2 violation produces oscillation,
and C4 drives the system towards coherence. The mechanism is identical; the scale
differs by roughly twelve orders of magnitude.
\end{remark}

\subsection{Synaptic plasticity as C4 at the neural level}

The brain does not maintain a fixed signed graph. The synaptic weights evolve
continuously in response to neural activity, a process called synaptic plasticity.
The Hebbian rule --- neurons that fire together wire together --- states that the
synaptic weight $w_{ij}$ between neurons $i$ and $j$ increases when they are
co-activated and decreases when they fire in opposition. In the language of signed
graphs, this is C4 operating on the edge signs of the neural graph in real time:
connections that reduce the coherence cost of the network are reinforced, while those
that increase it are weakened.

\emph{Interpretive remark:} Learning, in this picture, is not an additional
biological mechanism layered on top of physics. It is C4 applied to the neural
signed graph, exactly as C4 at the quantum level selects the coherence-minimising
configurations of $\Kfour$ and its composites.

The more precise formulation known as spike-timing-dependent plasticity (STDP)
reinforces a synapse when the presynaptic neuron fires before the postsynaptic one,
and weakens it in the reverse order. This directional temporal rule reflects the
half-cycle ordering of C1: which closure fires first determines the sign of the
coupling, exactly as the ordering of the two half-cycles of the $\Kfour$ pulsation
determines the sign structure of the coupled composite.

\subsection{The neural network as a composite closure}

The human brain contains approximately $10^{11}$ neurons and $10^{14}$ synaptic
connections. In the \PDL\ framework, this is a signed graph of very large size ---
a composite closure produced by many successive applications of $\Lop$, starting from
$\Kfour$ at the quantum level and passing through molecular, cellular, and tissue
levels. Its sign assignment evolves under C4 via synaptic plasticity. A coherent
state of this graph --- a configuration in which the fraction of frustrated triangles
is locally minimised --- corresponds, in this picture, to a stable thought or
cognitive state. Cognitive conflict is a local violation of C2 that generates
pressure towards resolution, exactly as a frustrated triangle in a signed graph
generates pressure under C4.

The self-representation hierarchy now has a natural neural correlate. $\Rop^1$ at
the neural level corresponds to the existence of a sub-network that encodes the
constraint structure of the global network --- in biological terms, a memory system
that models the organism's own coherence conditions (what works, what does not, what
maintains the closure). $\Rop^2$ corresponds to a sub-network that models the
modelling process itself: a system that represents its own representations. The
default mode network of neuroscience --- the brain's resting-state network, active
when the brain is not engaged with external stimuli and associated with
self-referential thought --- is the empirical candidate for this $\Rop^2$ sub-graph.

\emph{Interpretive remark:} The identification of consciousness with
$\Rop^2_{\mathrm{active}}$ and of the default mode network with the $\Rop^2$
sub-graph is coherent with the \PDL\ formalism and with the neuroscientific evidence,
but it is not derived from C1--C4 as a theorem. It is the target of future work in
the DL series, beginning with DL04.

\subsection{From grain of sand to amoeba: the threshold made concrete}

The distinction between a non-living composite closure and a living one can now be
stated precisely in terms of the framework developed above.

A grain of sand satisfies C1--C4 at every level from the quantum to the mineral: its
atoms pulse, its ionic bonds minimise coherence cost, its crystal structure is
irreducible (C3). But it does not satisfy $\Rop^1_{\mathrm{active}}$. It has no
sub-structure coupled via $(A)\wedge(B)$ to a production mechanism. It does not
generate $\Gamma_{i+1}$ from $\Gamma_i$. It persists passively as long as its
environment permits, and fragments when external forces exceed its coherence cost.
C4 has selected its configuration as optimal at its level. It is blocked there.

An amoeba also satisfies C1--C4 at every level from the quantum to the cellular.
But it additionally satisfies $\Rop^1_{\mathrm{active}}$: it contains a sub-structure
(its genetic material) that encodes $\Sop(\Gamma)$ via a sign-preserving map, and
this sub-structure is coupled via $(A)\wedge(B)$ to the production of new closures
of the same type. It generates $\Gamma_{i+1}$ from $\Gamma_i$ with controlled
variation $\delta \approx \delta^{*}$. It dies --- releasing the resources and
constraints that allow $\Gamma_{i+1}$ to exist --- and it evolves, because C4 at
the lineage scale selects the lineages of lowest average coherence cost, not the
individuals of lowest instantaneous cost.

The boundary between grain of sand and amoeba is not a boundary of substance. It
is a boundary of level and of recursive structure. Both are made of the same atoms,
governed by the same four axioms. The difference is the number of times $\Lop$ has
been applied, and whether the result satisfies $\Rop^1_{\mathrm{active}}$.

% ============================================================
\section{The Thermodynamic Ceiling and Combinatorial Entropy}
\label{sec:entropy}
% ============================================================

The rarity law established in Computational Result~\ref{thm:rarity} has a direct
thermodynamic interpretation. The fraction $f(n)$ of coherent configurations
decreases doubly exponentially with $n$: $f(3) = 1/2$, $f(4) = 1/8$, $f(5) =
1/64$, $f(6) = 1/1024$. At level $n$, C4 must select a coherent configuration from
a space of $2^{\binom{n}{2}}$ possibilities, of which only $f(n) \cdot
2^{\binom{n}{2}} = 2^{n-1}$ are coherent. The \emph{combinatorial entropy cost} of
maintaining coherence at level $n$ is proportional to $-\log_2 f(n) = \binom{n}{2}
- n + 1$, which grows as $n^2/2$ for large $n$.

On the other side of the balance, the gain from each additional level of
self-representation diminishes: the first level of $\Rop$ provides the largest
reduction in lineage coherence cost, and each subsequent level provides a smaller
marginal gain. The ceiling $n^{*}_{\mathrm{max}}$ is the level at which the
combinatorial entropy cost exceeds the marginal self-representation gain.

This analysis is the PDL-theoretic counterpart of Prigogine's dissipative structures
\cite{Prigogine1977}: the analogy is structural, not metaphorical. Prigogine showed
that systems far from thermodynamic equilibrium can spontaneously organise into
structures of increasing complexity, bounded by the capacity of the environment to
absorb entropy. The \PDL\ framework identifies the mechanism driving this
organisation (C4 applied to composite closures) and the mechanism imposing the
ceiling (the rarity law $f(n)$), both derived from the same four axioms.

The difference is one of depth. Prigogine's framework presupposes thermodynamics,
energy, and dissipation. The \PDL\ framework derives these from C1--C4.

% ============================================================
\section{Relation to Existing Literature}
\label{sec:literature}
% ============================================================

Several established programmes address the emergence of life and consciousness from
physical or informational principles. The \PDL-V conjecture differs from all of them
in a single structural respect: the axioms C1--C4 are prior to physics, not a
reformulation of physical law.

\emph{Integrated Information Theory} (IIT, Tononi \cite{Tononi2004}) proposes that
consciousness is identified with the integrated information $\Phi$ of a physical
system --- the information generated by the system above and beyond its parts. The
IIT postulates are stated at the phenomenological level and derive physical
requirements from them. The \PDL-V conjecture proceeds in the opposite direction:
from combinatorial axioms prior to experience, towards the conditions that force
self-representation. The condition $\Rop^2_{\mathrm{active}}$ has structural
similarities to high-$\Phi$ systems (both are irreducible to their parts, by C3 and
by the definition of $\Phi$ respectively), but the derivation is entirely different.
IIT has been criticised as unfalsifiable \cite{IIT2023}; the \PDL-V conjecture
provides explicit falsification criteria (Section~\ref{sec:PDLV}).

\emph{Autopoiesis} (Maturana and Varela \cite{Maturana1980}) defines life as
operational closure: a system that produces its own components and thereby maintains
its own organisation. This is structurally parallel to $\Rop^1_{\mathrm{active}}$ in
the \PDL\ sense: a closure whose encoding sub-graph drives the production of new
closures of the same type. The formal connection is not coincidental, but the
derivation differs: autopoiesis is a definition imposed from outside, while
$\Rop^1_{\mathrm{active}}$ is conjectured to be forced by C4.

\emph{Autocatalytic sets} (Kauffman \cite{Kauffman1993}) propose that life emerges
inevitably in sufficiently rich chemical spaces via mutually catalytic networks.
Kauffman's argument that autocatalysis is not improbable but inevitable beyond a
complexity threshold is structurally analogous to the claim that $n^{*}_{\mathrm{life}}$
is finite and forced by C4. The mechanism differs: Kauffman's argument is
combinatorial-probabilistic over chemical space, while the \PDL\ argument is
structural over signed graph hierarchies.

\emph{Logical depth} (Bennett \cite{Bennett1988}) measures the complexity of an
object by the computation time required to produce it from its minimal description.
Living organisms have high logical depth. The level $n$ in the \PDL\ hierarchy is a
structural analogue of logical depth: a closure at level $n$ requires $n$ applications
of $\Lop$ to construct from $\Kfour$.

The decisive difference between all these frameworks and the \PDL-V programme is the
following. Every existing framework either presupposes physics, chemistry, or
phenomenal experience as its starting point. The \PDL\ axioms C1--C4 presuppose
none of these. Physics is the first theorem of C1--C4. Life, in the \PDL-V conjecture,
is a later theorem in the same chain.

% ============================================================
\section{Epistemic Status Summary and Research Programme}
\label{sec:programme}
% ============================================================

\subsection{Complete epistemic status table}

\begin{center}
\begin{longtable}{p{6.5cm}p{2.8cm}p{4.5cm}}
\toprule
\textbf{Claim} & \textbf{Status} & \textbf{Reference} \\
\midrule
\endhead
$\Sop(\Kfour) \cong \Kfour$ unconditionally & Comp.\ result & Script~1, 8/8 \\
$\Sop(K_n) \not\cong K_n$ for $n \geq 5$ & Comp.\ result & Script~1, 0/16, 0/32 \\
$\Rop^1_{\mathrm{strong}}$ in 4/8 $\Kfour$ configs & Comp.\ result & Scripts 1, 4 \\
$\Rop^1_{\mathrm{weak}}$ forced for $n \geq 4$ & Comp.\ result & Script~5, 8--64/8--64 \\
$\delta^{*} = 1/2$ & Comp.\ result & Script~3 \\
Rarity law $f(n) = 2^{n+1-\binom{n}{2}}$ & Comp.\ result & Script~3 \\
\midrule
Scale-invariance of C1--C4 & Conj.\ (struct.) & Section~\ref{sec:scaleinvariance} \\
Existence of $n^{*}_{\mathrm{life}}$ (finite) & Conj.\ (struct.) & PDL-V Part~1 \\
C4 selects $\Rop^1_{\mathrm{active}}$ & Conj.\ (struct.) & PDL-V Part~2 \\
Death forced by C4 at lineage scale & Conj.\ (struct.) & PDL-V Part~3 \\
$n^{*}_{\mathrm{max}}$ finite & Conj.\ (struct.) & PDL-V Part~4 \\
\midrule
Synapse as signed edge & Interp.\ remark & Section~\ref{sec:synapse} \\
STDP as C1 temporal structure & Interp.\ remark & Section~\ref{sec:synapse} \\
Consciousness as $\Rop^2_{\mathrm{active}}$ & Interp.\ remark & Section~\ref{sec:synapse} \\
Default mode network as $\Rop^2$ sub-graph & Interp.\ remark & Section~\ref{sec:synapse} \\
\midrule
Value of $n^{*}_{\mathrm{life}}$ & Open problem & DL02 target \\
Formal derivation of C4 cost advantage & Open problem & DL03 target \\
$\Rop^2$ threshold and consciousness & Open problem & DL series \\
\bottomrule
\end{longtable}
\end{center}

\subsection{Planned DL document series}

\begin{description}[leftmargin=2em]
\item[DL01 (this document)] Foundations of the \PDL-V conjecture; lifting operator
$\Lop$; constraint structure $\Sop$; self-representation $\Rop^k$; computational
verification on small cases; epistemic stratification.

\item[DL02] Formal derivation of the C4 cost advantage for $\Rop^1_{\mathrm{active}}$
closures over non-representing ones; first estimate of $n^{*}_{\mathrm{life}}$ from
the rarity law and the marginal gain function.

\item[DL03] The birth--death cycle as a theorem: proof that immortal closures have
strictly higher lineage coherence cost than finite-duration lineages with $\delta =
\delta^{*}$; derivation of $\delta^{*}$ from C1--C4 at arbitrary level.

\item[DL04] The synaptic mechanism as a formal application of $\Lop$ at the neural
level: derivation of the Hebbian rule and STDP from C4 applied to the neural signed
graph; first formal steps towards $\Rop^2_{\mathrm{active}}$.

\item[DL05 and beyond] The $\Rop^2$ threshold; consciousness as a theorem; the
ceiling $n^{*}_{\mathrm{max}}$ as a computable integer.
\end{description}

% ============================================================
\section{Conclusion}
\label{sec:conclusion}
% ============================================================

The \PDL\ framework has established, through documents D01--D53, that the fundamental
constants and dynamical laws of physics are derivable from four combinatorial axioms
on finite signed graphs, without presupposing spacetime, particles, or fields. The
DL programme asks whether the same axioms, applied hierarchically via the lifting
operator $\Lop$, force the emergence of life and consciousness as structural
necessities rather than as contingent accidents.

The present document establishes the foundational objects of this programme. The
constraint-structure operator $\Sop$ is defined and shown to have $\Kfour$ as its
unique fixed point among complete graphs. Weak self-representation
$\Rop^1_{\mathrm{weak}}$ is shown computationally to be universally present for all
coherent closures with $n \geq 4$ vertices. The optimal inter-generational variation
$\delta^{*} = 1/2$ is identified as the half-pulsation of $\Kfour$. The rarity law
$f(n) = 2^{n+1-\binom{n}{2}}$ establishes the thermodynamic ceiling $n^{*}_{\mathrm{max}}$
as a combinatorial necessity.

The \PDL-V conjecture is not yet a theorem. What it is, is a precisely formulated,
computationally grounded, falsifiable research programme that asks the deepest
question one can ask of a foundational framework: if it truly precedes physics, does
it also precede life?

The answer, if it comes, will come through the same discipline that produced D01
through D53: one definition at a time, one script at a time, one proof at a time.

\clearpage

% ============================================================
\bibliographystyle{unsrt}
\bibliography{DL01_references}
% ============================================================

\end{document}